% TRANSLATION-PRODUCER DRAFT — UNCHECKED
% Work: Emmy Noether, Paper 5, Rationale Funktionenkörper — Korean U04
% Current pointer: NOETH-DE-AUTH-v003-20260804; pointer SHA-256 932FEDC1735A41A9CF71D15A6C662A468A4CAD016AE8B3DECDF9A71E8BA7F197
% Current German authority: C:\Users\Floris\Documents\interlanguage\03_projects\noether\07_german_canon_control\candidates\NOETH-DE-ED-0001\Noether_German_NOETH-DE-ED-0001.tex
% Authority identity: 2,153,565 bytes; SHA-256 D1F06B311F6CBD991DD247D745DD9A72DDE326A20396DF43CFE0C8EDB1593CDB
% Paper 5 body interval: whole-source lines 4535--4573; 8,185 LF-normalized UTF-8 bytes; SHA-256 99BD68A8DBD9861EFF0CDBE26CB365C3306EDF15BD93A6C4C10B9F25419D5CAE
% Unit source: local lines 31--38 / whole-source lines 4565--4572; 1,770 LF-normalized UTF-8 bytes; SHA-256 FF5058094D557ECF29D2EA4A37762067EF5936EAD529E70AF5C5FA0E6B063230
% Producer boundary: Korean translation only. No source/scan comparison, Korean or formula review, compilation, rendering, assembly, packaging, certification, approval, German adjudication/edit, or self-check was performed.
% Authority identities are routing metadata, not a producer source audit or approval.
% Handoff state: UNCHECKED; all Korean wording and preserved TeX require an independent Korean checker.

현재로서는 정수성 기저가 존재하는 체의 한 부류만을 제시할 수 있다. 즉 그 체가 $n$개의 다항식으로 이루어진 한 계를 포함하고, 이 다항식들의 최고차항들의 동차 종결식이 0이 아니면 정수성 기저가 존재한다. 이 기저는 다음 선형형식에 대하여 그 체에서 불가약인 방정식의 모든 $u$에 관한 계수들로 주어진다:
\[
  u_0=u_1x_1+u_2x_2+\cdots+u_nx_n.\srcfnmark{2)}
\]
\srcfntext{2)}{부정원이 하나인 경우, 이 불가약 방정식은 Lüroth 정리를 증명하기 위한 E. Steinitz의 논의에 나온다.}
특히 종결식의 값이 계수영역의 단위원과 같으면 표현의 정수성도 보장된다. 이 체의 부류에 속하는 한 예는 라그랑주 종영역이다. 기본대칭함수들의 종결식 값이 $\pm 1$이기 때문이다. 또 다른 예는 정수성은 없지만, 대수적으로 독립인 다항식을 단 하나만 포함하는 모든 체이다. 이때 정수성 기저는 모든 다항식을 포함하는 가장 작은 부분체의 Lüroth 함수와 동일하다. 이와 같이 $n$개 변수의 이차형식에 대한 모든 정칙 유리 사영불변식은 판별식의 정칙 함수라는 것이 알려져 있다.

위의 조건은 정수성 기저가 존재하기 위한 충분조건일 뿐, 결코 필요조건은 아니다. 임의의 $n$개 다항식의 종결식이 0이면서도 저 불가약 방정식의 계수들이 주는 정수성 기저를 갖는 체를 쉽게 구성할 수 있다. 따라서 가장 일반적인 경우에도 혹시 저 방정식의 계수들이 정수성 기저를 주는가라는 문제가 생긴다.
